\documentclass[11pt]{article}
%\usepackage{hyperref}
\usepackage{verbatim}
\usepackage{natbib}
\bibpunct[: ]{(}{)}{;}{a}{}{,}
\usepackage[margin=1.2in]{geometry}
\usepackage{setspace} 
\usepackage{graphicx}
\usepackage{url}
\usepackage{lscape}
\usepackage{subfig}

%\hypersetup {
%	pdfauthor={Omar Lizardo}, 
%	pdfpagemode = UseNone, 
%	pdfstartview=FitH, 
%	linkbordercolor = .7 .7 .8,
%	citebordercolor = .7 .7 .8,
%	}

\begin{document}

\title{Cultural Consumption, Culture-Based Talk, and Stratification:  Why `Cultural Matters' Matter\footnote{Previous versions of this paper were presented at colloquia at the University of Notre Dame and at the MIT/Harvard Economic Sociology Seminar series. I would like to thank Roberto Fernandez, Robert Fishman, Gene Halton, Mary Ellen Konieczny, Elena Obukhova and Ezra Zuckerman for helpful comments on earlier drafts.  Direct correspondence to Omar Lizardo, University of Notre Dame, Department of Sociology, 810 Flanner Hall, Notre Dame, IN, 46556  Tel: (574) 631-1855, Fax: (574) 631-9238, Email: \ttfamily{olizardo@nd.edu}.}}

\author{Omar Lizardo \\ \textit{University of Notre Dame}}


\date{\today}


\maketitle

\begin{center}
	Word count:  11,216 (tables and text)
\end{center}

\newpage
\begin{center}
	{\large Author's biography}
\end{center}
{\bf Omar Lizardo} is an Assistant Professor of sociology at the University of Notre Dame.  His research interests include the sociology of culture and taste, social networks, social theory and the sociology of globalization.

\newpage
\begin{center}
	{\large Cultural Consumption, Culture-Based Talk, and Stratification}
\end{center}
\begin{abstract}
\small I analyze survey data from a probability sample of Americans to  for the first time document the distribution and predictors of what I call ``culture-based talk'': essentially, the use of cultural content to generate conversation with friends and acquaintances.  I argue that this phenomenon---central to many \emph{theories} of the connection between culture and stratification---deserves more \emph{empirical} attention. I show that culture-based talk is a pervasive feature of the lives of most Americans, especially those who are culturally active. In addition, I show that rates of culture-based talk appear to be stratified by the usual markers of status (such as levels of schooling and generational location), and for some cultural domains (such as music), this stratification is independent of frequency of cultural engagement.  In this last case, the results show that members of high-status taste cultures are more likely to engage in culture-talk than members of low status taste cultures.  I close by outlining the theoretical and empirical implications of taking the phenomenon of culture-based talk seriously. 
\end{abstract}

\newpage
\doublespacing
\section{Introduction} 
Empirical research at the intersection of the sociology of culture and stratification has been primarily concerned  with exploring inequalities linked to social position in access to and engagement with symbolic goods  \citep{dimaggio00, vaneijck00, katzgerro04}.  Whether taste and culture consumption patterns---as constitutive of the lifestyles that mark Weberian status groups---continue to be associated with social position continues to be the central question in this literature \citep{tomlinson03, katzgerro04}.  This organizing concern can be appreciated whether the construct of ``social position'' is conceptualized in terms of educational attainment \citep{bourdieu_darbel91, bryson96, vaneijck99}, membership in broadly defined occupational groups \citep{bourdieu84, dimaggio_useem78a, dimaggio_useem78b, peterson_simkus93, tomlinson03}, aggregate measures of social class \citep{katzgerro_shavit98, katzgerro02, katzgerro06}, or constructed scalar indices of social status based on patterns of social intercourse across occupational groups \citep{chan_goldthorpe05, chan_goldthorpe07a, chan_goldthorpe07b}. 

The overall conclusion that has been reached in this still vibrant line of research is that lifestyle practices continue to be strongly shaped---albeit in increasingly complex ways---by social position \citep{holt98, katzgerro99, tomlinson03, emmison03, warde_etal08}.  This runs counter to the predictions of ``postmodernization'' and class ``destructuration'' theorists who foresaw the decline in the linkage between class and lifestyle in market-dominated ``consumer societies''\citep[e.g.][]{baudrillard06, kingston00, pakulski_waters96}.  The key insight is that generalized markers of social status and education (rather than social class narrowly defined) provide the best lens to understand cultural stratification patterns \citep{alderson_etal07,chan_goldthorpe07a, chan_goldthorpe07b, katzgerro02, katzgerro04, sullivan_katzgerro07, silva06, bennett_etal99, emmison03, ollivier04}.  Confirming the basic findings of an earlier generation of research \citep{baumol_bowen66, dimaggio_useem78a, ganzeboom82}.

I argue that the strong focus of the cultural stratification literature on the question of \emph{whether} lifestyle practices continue to be unevenly segmented according to class and status markers has resulted in a relative lack of attention to an equally important question:  what is the \emph{role} (or ``function'') of aesthetic consumption in the everyday lives of members of different social strata?  More precisely, given that members of certain social strata are more likely to engage culture and the arts more than other groups, to what \emph{uses} do they put the culture that the consume?  

This gap in the literature is surprising in light of the fact that the question of the role played by engagement with symbolic goods was an continues to be central to the theoretical tradition from which most analysts of the culture-stratification linkage depart in the first place \citep{collins75, dimaggio87, douglas_isherwood96, fine77}.  In this respect, Fine's basic complaint in relation to research of popular culture in the early 1970s, can be repeated vis a vis research on cultural stratification today.  This research, generally ignores the \emph{behavioral implications} implications of culture consumption, especially as they pertain to social interaction \citep[453]{fine77}. 

A key insight to be gained from these foundational sources is that if we cannot get a solid theoretical handle on the uses to which persons put the cultural knowledge acquired through aesthetic experiences then it is difficult to analytically justify the status of these symbolic resources as ``cultural capital'' in a given social setting \citep{mohr_dimaggio95, dimaggio96, dimaggio00, carter03}.  This is the reason why in their influential paper Lamont and Lareau \citeyearpar{lamont_lareau88} linked the notion of cultural capital to an explicit conceptual framework that highlighted the active deployment of a given repertoire of cultural tastes and lifestyle practices for the purposes of drawing symbolic boundaries.  DiMaggio \citeyearpar{dimaggio87, dimaggio00} is also clear---building on the early work of Collins \citeyearpar{collins75}---in proposing that aesthetic consumption of cultural goods must be connected to routine exploitation of cultural knowledge in social interaction within everyday settings (see  \citet{erickson96} and \citet{lizardo06a} for empirical applications and theoretical elaborations of this idea).    From this perspective, one substantively important way in which the set of cultural aptitudes that manifest themselves as cultural \emph{consumption} become actualized as cultural \emph{capital} is through the linkage of cultural engagement to everyday processes of informal socialization and thus status group reproduction in the Weberian sense \citep{dimaggio_mohr85, dimaggio96}.

\subsection{Analytical and empirical goals}
In this paper, I draw out the empirical implications of this set of starting ideas for our understanding of the \emph{social uses} of aesthetic consumption in interaction.  My analytical point of departure is the proposition that one of the primary ``uses'' of aesthetic consumption are social \citep{dimaggio87};  even if the ``gratifications'' that it provides can sometimes be personal and private  \citep{denora00, hennion01, holbrook_hirschman82}.  This implies that the everyday use of cultural knowledge and experiences gained through aesthetic consumption become a strategic---although naturally not the only---resource for the initiation social interaction in conversation.  This happens in both relatively formal \citep{long03} and informal \citep{fine77} settings.   I propose that inequalities in cultural access can be more profitably understood as reflecting inequalities in the ability and opportunity to engage in certain forms of informal interaction that Simmel \citeyearpar{simmel49} referred to by the name of ``sociability.''  

Empirically, this paper makes three main contributions.  First, I document the distribution of culture-based talk, as far as the data permit, in the American population.  The results indicate that culture-based talk is fairly widespread in the United States, with most Americans reporting having used culture from one of the five domains considered here at least fairly often in conversation with persons they know.  Second, I show that even if widespread, culture-talk is highly stratified by socio-demographic markers, in particular education (across both non-profit and commercial domains), age (primarily in commercial domains such as film and music) and gender (mainly when it comes to literature). Finally, I show that for domains in which status-based stratification based on talk is strong without corresponding levels of stratification based on consumption, culture-talk is stratified by membership in a given \emph{taste culture}.  Here, respondents more likely to belong to more artistically legitimated taste cultures engage in culture-talk more frequently.

Analytically, this research has important implications for our understanding of the link of culture and stratification in late-modern societies \citep{ollivier00}.  If the argument offered here is on the right track, then it becomes clear that for culturally-advantaged members of high-status groups the deployment of cultural capital \emph{in interaction} \citep{eliasoph_lichterman03} during face-to-face talk serves as an key \emph{relational mechanism} \citep{tilly98} that links social position to aesthetic consumption.  This mechanism explains both status-based advantages premised on the maintenance of (homo)social bonds, as argued in the classic Weberian tradition, as well as those social advantages tied to the ability to connect across distant social positions in the social structure \citep{erickson96, lizardo06a, lin99, pachuki_breiger10}.   Rethinking informal socialization as sociability allows us to reconsider the role of cultural content in the constitution of social relations.  In particular, I conclude that there might not be any necessary relationship between the ``importance'' of the discussion-content and the importance of the social tie in question (and in many cases this relationship will be reversed).

\section{Culture as a social resource:  Specifying the phenomenon}
The empirical issue that is the focus of this paper pertains to the ``behavioral implications'' of the aesthetic consumption of popular (and fine arts) culture \citep{fine77}.  This phenomenon can best be characterized through the frequently noted observation that cultural knowledge obtained from aesthetic consumption experiences will become a likely content used for the initiation and maintenance of conversational interaction \citep{dimaggio87, collins75}.  I refer to this (presumably) pervasive phenomenon as \emph{culture-talk}.  Culture-talk can be defined as the deployment of cultural knowledge obtained from aesthetic consumption experiences as a resource to generate conversation.  In a nutshell:  the discussion of ``cultural matters'' with friends and acquaintances.

Culture-talk is one of the most common facilitators of \emph{sociability} in late-modern societies.  Sociability is social interaction that is not driven by explicit instrumental purposes, but that is keyed to the continuing repetition of episodes of social intercourse for their own sake \citep{simmel49}. This is a form of social intercourse explicitly dissociated from direct instrumental pursuits, as a \emph{form} of association ``which does not have a strictly economic or business purpose'' but access to which (and restrictions on) are a key marker of the status group boundaries \citep[114]{weber94}. It is in this respect that sociability can be thought of as the ``play form'' of the more general process of association (coming together from some purpose).  This type of interaction would be difficult to successfully realize by relying on traditional ``instrumental'' contents (e.g. work, politics, etc.), but requires instead ``ludic'' content that serves to simply sustain the focus of attention in the interaction itself \citep{collins04}.  I argue that the ``entertainment talk'' typical of those who are active culture consumers is one of the forms of cultural knowledge that is most apt for these ``playful'' purposes.

This Simmelian mechanism is at the center of Collins' microinteractionist theory of culture and stratification \citep{collins75, collins81, collins00, collins04}.  Simmelian sociability is primarily built around what he refers to as ``conversation for enjoyment.''  These types of conversational rituals are important because they signal a removal from instrumental realms laden with power-dynamics,  status differences and ``arms-length'' relationships built for explicit, narrow purposes.  Instead, conversation for enjoyment signals that we have entered  into a homo-social realm based on (situational) status equality and the privilege of status inclusion \citep{milner04}.  For this reason, it is to be expected  that culture-based talk should be more likely to happen among friends than among strangers \citep[124]{collins75}.  Nevertheless, certain forms of culture-talk---especially those involving cultural content presumed to be spread evenly across groups---provide the means to situationally relax status groups boundaries by providing a baseline form of cultural knowledge that can allow otherwise socially distant persons to ``break the ice'' \citep{erickson96}.  

As I will elaborate more fully  below, this has important implications for the way in which researchers link cultural content with the subjective and functional significance of social ties.  Traditional methodological practice in network analysis is governed by the presumption that ``important'' relations can be induced by inquiring about contacts with whom ``important'' cultural matters are discussed \citep{burt84, marsden87}.   This presupposes that the importance of the contact (in a functional or sociological sense) matches the substantive ``importance'' of the cultural content through which the relation is defined. Following the Simmelian line of reasoning outlined so far (which suggests that ludic content is required for enjoyable sociability), I suggest that it makes more substantive and phenomenological sense to inquire about seemingly ``unimportant'' cultural matters.  This is particularly important if the goal is get at those social relationships that are most meaningful for the person and of most theoretical relevance for the analysis \citep{pachuki_breiger10, mische10}.

\section{Theoretical and Substantive Importance of Culture-Based Talk}  
Why should the phenomenon of culture mediated interaction be of interest to sociologists interested in the linkage between culture and stratification?  While the connection between aesthetic consumption and social interaction (and thus the uses of culture to build and maintain social networks) has been a relatively neglected issue in the current sociology of culture literature the phenomenon has been clearly perceived to be a central issue in some of the foundational theoretical statements in the field. This literature provides us with some leads as to why we should find that culture-based talk is a relatively generalized phenomenon and why it should be stratified by the familiar markers of status.  

\subsection{Relationship formation}
Cultural knowledge gained through aesthetic consumption is a key resource with which to form \emph{new} social relationships.  ``culture-talk'' is especially useful at the \emph{initial} stages of ``the acquaintance process'' \citep{newcomb61}.  This is a ``need'' that is structurally induced by the fact that, as recent research on the dynamics of personal networks show, new relationships are constantly being formed and old ones renewed (or dropped) during the adult life-cycle \citep{bidart_degenne05, burt00}:  The most surprising thing about ``intimacy'' and ``friendship'' is how unstable it tends to be \citep{wellman_etal97}. Culture-talk should be intimately implicated in the renewal and formation of social relationships.

Culture-talk is a lubricant that keeps conversation going when more intimate content is not yet available (or is situationally or institutionally precluded).  This allows persons to search for ``deeper'' areas of commonality and agreement without much risk (as would inquiring about politics or economic position), so that the interaction can proceed smoothly.   When individuals do not have a previous social or cultural bases for interaction based on common membership in the same social network or common biographical experiences they instead can rely on uncontroversial cultural topics to initiate and sustain interaction \citep[460]{fine77}. In this way, culture-talk is able to facilitate the type of social interaction that produces intimate relationships out of initially weak bonds of acquaintanceship \citep[441]{dimaggio87}.  

\subsection{Status-based solidarity in small groups}
Status-based exclusion occurs largely through micro-mechanisms governing entry (or determining systematic exclusion from) class-structured micro-worlds organized around specific conversational interactions.  Culture-based talk increases the cohesion and social closure of small groups, providing a way to segment them from the larger social structure \citep{fine77}.  For instance,  Simmelian urban social circles are produced and reproduced through shared tastes and practices---e.g. concert going, visiting the museum, etc.---intrinsically connected to culture-based talk \citep[797]{kadushin66}. 

Because is a relational mechanism for ritual identification, the construction of social relationships and the formation of networks of trust and sociability, networks formed through shared cultural content are potentially available as social capital, facilitating ``the attainment of such social rewards as desirable spouses and prestigious jobs'' \citep[443]{dimaggio87}.  One functional consequence of shared aesthetic consumption experiences is to serve as the non-instrumental content that justifies the existence of status-segregated conversational rituals in the first place \citep{ostrower98}. As \citet[125]{collins75} notes, ``\ldots [i]n some social circles{\ldots}entertainment talk makes up most of the content of conversations; in all classes, it is important.'' 

\subsection{Linkages within loosely bounded and sparsely knit networks}
Culture-based talk is not only  the ``bond'' that keeps status groups together (thus sustaining discontinuities in the social structure), but also the ``solvent'' that allows for sporadic connections across status lines.  Culture-based talk provides new forms of high-status group solidarity not built around traditional notions of commonality anchored to geography, community and shared institutional affiliations and kinship structures \citep{holt98}. This process acquires more importance as these network ties have been transformed in the contemporary context of increasing geographic mobility \citep{griswold_wright04} and the decline of the traditional bonds characteristic of traditional local communities \citep{wellman79}, which were centered on kinship and spatial contiguity \citep{dimaggio_mohr85}.  As the personal networks of socially mobile members of the upper-middle class have become more fluid, loosely knit, more discontinuous and less tied to propinquity and kinship \citep{meyrowitz85, wellman_wortley90, wellman_etal88}, the arts and popular culture have come to increasingly serve as the ``default'' forms of knowledge that connect people across distinct interaction settings \citep{dimaggio87}.

The symbolic meaning  of popular culture can be safely taken to be shared by most persons, which makes them an appropriate topic of conversation even when persons are not well acquainted  \citep[460-461]{fine77}. This form of culture-talk may serve to connect persons across different positions in the social structure and across segregated taste publics \cite{ikegami05, mische10}.   This allows for the formation of social connections---even if sometimes brief and sporadic, as in so-called ``weak ties''---when there exist large socio-demographic differences \citep{lizardo06a}.  The deployment of this type of cultural knowledge in interaction in more instrumental settings may even facilitate coordination across hierarchical positions \citep{erickson96}. This can happen even as the deployment of this type of cultural knowledge simultaneously serve as its own basis for exclusion (as when men of different classes can bond over ``sports talk'' that excludes women and those men not familiar with the local sports culture \citep{erickson96, purcell07, pachuki_breiger10}).  

In sum, some of the most important lines of theory in the culture and stratification literature point to the central importance of the phenomenon that we have labeled and characterize as \emph{culture mediated social interaction}.  These theoretical streams point to this phenomenon as a pervasive, ubiquitous feature of late modern societies.  According to this line of theory, the phenomenon of culture-based talk is key to understanding the dynamic, interactional origins of the social discontinuities and status-based social boundaries that characterize the modern social order \citep{illouz97}. This makes its almost total neglect in the more empirically oriented recent literature in culture and stratification an even more puzzling oversight. 

From this perspective, the formal and informal system of production of aesthetic goods provide as ``input'' the very content that serves as the scaffolding on which everyday episodes of ``talk'' and sociability are built.  The socially segmented micro-orders based on face-to-face (and increasingly electronically mediated) interaction require a constant stream of new interactional content in order to be sustained and thus reproduce themselves over time \citep{collins81, collins04, dimaggio87, lizardo06a}. This implies that the stratification of culture-based talk should mirror in non-trivial respects patterns of stratification found in the consumption of aesthetic goods and experiences.

\subsection{Research Questions}
The lines of theory that put culture-based talk front and center in processes of cultural stratification are wide ranging.  In the empirical analysis that follows I cannot explore the full range of implications of these theoretical proposals.  Because so little is known about the distribution and correlates of culture-mediated talk, I will limit myself to the following research questions:

\begin{enumerate}
\item \emph{What is the relative distribution of the frequency of culture-based talk in the contemporary United States}? Insofar as culture-based talk is a key process in current theories of the culture/stratification link, we should expect it to a fairly widespread phenomenon.  Because there has been virtually no quantitative research on the issue, it is hard to estimate with certainty how much is ``widespread.''  Given what we know from the few qualitative studies that have been conducted, we can  expect to observe close to a majority of the population reporting high or moderately high frequencies of culture-based talk.\footnote{As with any study that uses arts participation data in order to draw larger conclusions about the operation of culture in society at large, the big caveat is of course that any observed frequency of culture-mediate interaction is probably an underestimate, as persons can only report interacting with friends on the basis of the cultural forms included in the survey.  Clearly, persons who are observed to not interact with friends very often on the basis of content taken from the cultural domains included in the survey, may still be doing so on the basis of cultural forms not included in the survey.}

\item \emph{Does the propensity of culture-based talk differ across socio-demographic categories?} While most persons are indeed likely to engage in some form culture-based talk, it is to be expected that as a general rule, culture-based talk should be stratified along indicators of status and prestige as well as such important status markers as age.  We should expect this to be the case for two reasons:  first, culture consumption is itself stratified, with more educated individuals being more likely to engage in a wide range of aesthetic consumption activities than the less educated \citep{bourdieu84, bourdieu_darbel91, vaneijck99} insofar as culture-based talk is much more likely to be observed among active culture consumers than among inactive ones, then there should be a parallel between the predictors of consumption and the predictors of culture-based talk.  

Second, opportunities for social interaction and such indicators of social connectivity as network size and proportion non-kin in the typical discussion networks are stratified by the same types of status and prestige indicators, in particular educational attainment \citep{marsden87, fischer82, burt90}, age, gender and race.  Insofar as culture mediated interaction is more likely to take place among members of informal cliques, we should expect that those persons most likely to have wide-spanning, loosely knit networks of acquaintance, composed of elective relationships of ``friendship'' (in contrast to ascriptive relations of kinship) should be more likely to engage in culture-based talk.  Most research on personal discussion networks indicates that the higher educated and those in mid-adult life are more likely to enjoy these types of local relational environments than the less educated or the elderly \citep{burt90, marsden87, erickson96, lizardo06a, fischer82}. 

\item \emph{Does the propensity of culture-based talk differ across taste cultures?} Theory and research in culture and stratification suggests that this is likely to be the case.  A key tenet of contemporary cultural theory is that audiences are not only segmented by socio-demographic markers such as levels of schooling or occupation, but that they are also partitioned according to \emph{taste cultures} \citep{gans99, peterson_dimaggio75}.  The basic hypothesis is that there exist systematic differences in the way in which members of different taste cultures appropriate cultural goods.  These distinct styles of engagement might be relevant for predicting taste-based differences in the relatively likelihood that culture will be a subject of conversation.  For instance, culture-talk among respondents who are attracted to traditionally ``artistic'' domains is likely to be based on those qualities of aesthetic goods that have been the subject of standardization and institutionalization in school curricula or likely to be explicitly formulated and elaborated within informal networks of enthusiasts and cultural professionals linked to critics, journalists and other ``genre-making'' gatekeepers \citep{bourdieu93a, benzecry09, frith98, weber77}.\footnote{The extent to which the aesthetic qualities of a given cultural form lend themselves to verbal expression is of course not inherent in the form itself, but dependent on the institutional history of the cultural domain in question \citep{bourdieu93a, frith98, regev94}. } 

In particular, cultural orientations consistent with the institutionalized demands of highly elaborated fields of cultural production---what Bourdieu \citeyearpar{bourdieu84} referred to as ``the aesthetic disposition''---are more likely to generate culture-talk in the context of informal networks of acquaintanceship and solidarity within certain strata.  Drawing on recent theory and research on the phenomenon of \emph{intellectualization} within artworlds \citep{baumann01,ferguson98}, I hypothesize that the presence (or relative absence) of an already clearly specified, highly institutionalized critical vocabulary capable of being recruited for culture-based talk, can be used to explain why we may observe categorical inequality in the underlying propensity to engage in culture-based talk \emph{in the absence} of a correlative sort of differentiation in terms of frequency of cultural engagement in that domain.  
\end{enumerate}

\section{Data and Variables}
\subsection{Data Source}
In what follows I examine data obtained from the Arts and Culture Supplement to the National Social Survey fielded in 2002 ($N=1,002$) and 2004 ($N=801$) \citep{NSS2002, NSS2004}.  The NSS is a weighted (by age, education, marital status, gender, race, ethnicity) probability sample respondents collected through computer-aided telephone interviews.   These surveys is particularly appropriate for our purposes because they contains information not only about culture consumption experiences, but also about the likelihood that individuals engage in informal social interaction featuring cultural content taken from consumption in the same domains.  Because the items of interest were fielded in both waves, I pool the samples for which I have complete data on all socio-demographic variables from the 2002 ($N=955$) and 2004 ($N=776$), which leaves us with a total combined sample of 1,731 respondents.

\subsection{Culture-Mediated Social Interaction Survey Items}
As part of the survey, respondents were asked the following question:

During the past twelve months, approximately how often have you had a discussion with someone you know about any of the following:
\begin{enumerate}
\item  \emph{A book, a poem, or a story?}
\item  \emph{A musical performance, live or recorded.}
\item  \emph{A movie?}
\item  \emph{A play or other dramatic performance?}
\item  \emph{An art show or work of art?}
\end{enumerate}

Respondents were asked to choose from the following response categories: 1) very often, 2) fairly often, 3) not very often, 4) or not at all. 
 
\subsection{Aesthetic Consumption Survey Items}
Respondents were also asked to report the frequency with which they engage in various forms of aesthetic consumption:\footnote{How does the NSS compare to other Arts participation Surveys?  When it comes to book reading, the 2002 NSS estimates that about 55.7\% of respondents report reading books ``Fairly Often'' or `` Very Often.''  The corresponding estimate for the 2002 SPPA (Survey for Public Participation in the Arts) is about 58\%.  When it comes to going to an art museum, about 12\% of NSS respondents report going ``Fairly Often'' and about 3.7\% of respondents report going ``Very Often''; these correspond well with SPPA 2002 data based on more specific questions on attendance frequencies:  13.4\% of SPPA respondents report having gone to the museum between 2 and 5 times in the last year, and 3.2\% report having gone more than five times.  In the NSS, about 9\% of respondents report going to a live performance ``Very Often.''  In the 2002 SPPA about 7.3\% of respondents report having gone out to a Jazz, Classical, Opera, Musical, Ballet or Dance performance between 2 and 5 times in the last year.  In all, the NSS appears to do a good job of capturing rates of arts participation that are comparable to that obtained from other National Surveys.}
 
\begin{enumerate}
\item  \emph{Books}:  Approximately how often during the past 12 months, have you read a book?
\item  \emph{Movies}:  Approximately how often during the past 12 months, have you been to the movies? 
\item\emph{Music}: Approximately how often during the past twelve months would you say you listened to music at home, for instance on CDs, tapes, records, a computer, or on the radio?
\item  \emph{Live performances}: Approximately how often did you go to any live performances such as plays, concerts, musicals, or dance?
\item  \emph{Art}:  Approximately how often did you go to any art shows or museums during the past 12 months? 
\end{enumerate}  

\subsection{Data Limitations}
Before proceeding with the analysis, it is important to note some important limitations of the data at hand.  First, the most salient limitation of the data source that I rely on in this study concerns the limited range of cultural domains included in the NSS for which information on the rate of culture-based talk is available.  In particular, as befits an arts participation survey, questions about culture-based talk are weighted toward the most legitimate forms of aesthetic practice but are silent when it comes to more popular culture consumption experiences \citep{dimaggio_mukhtar04}. In particular, past research suggests that culture-based talk centered around sports \citep{erickson96}; Fashion and Cuisine \citep{dimaggio87} and in particular television \citep{fiske87} are a consistent generator of culture-based talk in formal  and informal settings (e.g. the workplace).  Thus, we can be certain that the rates of culture based talk reported below \emph{underestimate} by a substantial amount its prevalence in the American population.   The absence of information on culture-based talk in more popular domains may also lead to an \emph{overstatement} of its connection to markers of stratification, such as educational credentials; although it is reasonable to suppose that we may be likely to observe education-based stratification in rates of culture-based talk in the absence of education-based stratification in consumption rates.  

Another important issue concerns the fact that there is no complete correspondence between the culture-talk and culture-consumption questions.  This becomes relevant whenever we attempt to compare the relative salience of socio-demographic differences across the talk and consumption domains.  This problem is most salient in the case of the ``live performance'' items.  While the culture-talk item deals with the more restricted domain of live drama, the consumption question deals with the broader domain of all live performances, including musicals and dance.  This obviously introduces error into any attempt to compare the relative importance of status markers across the two domains.  In the case of Movies and the Arts, the consumption questions are more restrictive than the talk questions, since the former deal explicitly with engagement with film and art outside of the home (going to the movies, visiting a gallery or museum), while the latter deal with the use of film or arts culture in general as a conversational topic.  It is of course likely that many respondents engage film and the arts without having to go the movie theater and the gallery, and this cultural knowledge should still be available for interaction all the same.  While we should expect broad positive correlations between more restricted and more generalized forms of participation \citep{dimaggio_useem78a, sullivan_katzgerro07}, the lack of strict congruence in the live performance domain should be kept in mind when interpreting the results. 

One final limitation of the data source is the fact that we have to rely on a single set of survey items to measure culture-based talk.  Survey questions have the advantage of allowing us to draw more general conclusions and subject the data to a powerful set of statistical manipulations.  However, they have the disadvantage of glossing the \emph{context} of consumption as a mediator of the extent to which given aesthetic consumption experiences will generate culture-based talk. For instance, some genres have built-in rituals of audience reception in which ``interaction'' is expected to occur \emph{within} the consumption site and \emph{during} the performance (e.g. a rock concert) and which thus carry little expectations that the cultural-knowledge gained therein will be transposed to other informal contexts as a generator of conversation.  For instance, culture that is consumed according to socially constructed standards of audience ``propriety'' (e.g. sitting quietly in rapt attention) on the other hand, may in turn generate a lot of culture-based talk precisely because direct interaction among audience members (verbal and non-verbal) is frowned upon during the reception process.\footnote{Cultural domains that exhibit wide variation in the context (and thus the expected audience conventions) of reception, may be seen to have a ``bi-modal'' distribution of their likelihood of being used for culture-based talk, with non-negligible proportions of active consumers reporting not using (as well as using) culture for that domain as a conversational resource. }  

\section{Analysis and Results}
\subsection{Distribution of culture-based talk in the United States}
\subsubsection{Overall frequency}
Figure~\ref{fig1} shows the univariate distribution of culture-based talk across the five broad realms of cultural activity included in the 2002 and 2004 National Social Survey.\footnote{the circle is the observed percentage and the vertical line represents the 95\% confidence interval, calculated using the ``exact'' binomial estimate for the observed proportion.}  The distribution of culture-based talk frequencies across all five cultural domains supports the expectation that persons use culture as a discussion topic in a fairly regular manner.  78\% of respondents report frequently (either ``fairly'' or ``very'' often) using \emph{at least one} of the cultural domains included in the survey as a subject matter in conversation with friends.

\begin{center}
	Figure~\ref{fig1} about here
\end{center}
 
When it comes to the three most popular cultural domains  (movies, literature and music), it is clear that a majority (or close to a majority) of respondents report having discussed matters related to that cultural domain with a friend.  About 60\% of persons report talking to friends about movies ``very often'' or ``fairly often''; the corresponding percentages for books and music are 52\% and 46\%. Even for the two cultural domains which are expected to have the most exclusive audiences (in terms of such socio-demographic characteristics as education and age), in this case theater and the arts, 26\% and 19\% of respondents report relative frequent rates of culture-based talk respectively. 
 
\subsubsection{Frequency among culturally active respondents}
The raw estimates of culture-based talk give a somewhat misleading picture of the extent to which the phenomenon occurs. It is reasonable to suppose that a necessary precondition for persons to deploy any type of cultural knowledge during interaction is that they have acquired this knowledge.  It is a good bet that knowledge about symbolic goods is acquired through direct engagement with the formal and informal culture industries (and non-profit organizational fields) that deliver these goods.  This means that we should expect the relative chances of observing high rates of culture-based talk to be mediated by the extent to which the persons actively engages the cultural field in question.
 
Figure~\ref{fig2} shows the conditional univariate distribution of culture-based talk among those respondents who report engaging in the corresponding cultural activity either ``fairly often'' or ``very often.''  I find that not surprisingly, rates of culture-based talk among active consumers of the cultural forms in question are substantially higher than the average.  About 90\% movie, 71\% of literature and 51\% of music consumers report engaging in relative frequent rates of culture mediated interaction.  The most dramatic relative increases in relation to the average is observed for the two cultural domains associated with the non-profit ``fine arts'' field:  53\% of active dramatic theater consumers and 61\% of active art consumers report discussing this type of cultural content with a friend.  

\begin{center}
	Figure~\ref{fig2} about here
\end{center}
 
Interestingly, the univariate distribution for culture-based talk based on music shifts the least after consumption is taken into account.  This indicates that this type of cultural content is deployed routinely without necessarily being tightly connected to consumption.  It can also mean that the ``uses'' of music as cultural form may be unique:  music may constitute both a ``social'' resource for the cementing of connectivity in social relationships, and also a private ``resource'' for personal enjoyment and the customization of ``life-spaces'' \citep{denora00}.
 
Another noticeable fact is that both movies and literature seem to be ``inherently social'' cultural domains in these data.  It appears to be almost (statistically at least) ``non-sensical'' to be an active movie consumer or book reader without using this cultural knowledge in interaction.  The concrete mechanisms governing the way in which these two cultural forms are deployed in everyday settings may be distinct however. The fact that in these data the proportion of respondents who actively engage the cinematic field \emph{without} at least occasionally using this cultural knowledge in conversation corresponds well with Fine's \citeyearpar{fine77} original proposal that cultural goods with widespread appeal would also be the most likely to be mobilized in interaction in small groups. Research on book groups by Long \citeyearpar{long03} and Illouz's \citeyearpar{illouz03} research on Oprah's televisual ``book club'' suggest that books are a primary conversational resource among a wide variety of audiences.  However, literature conversations tend to occur within the relatively formalized ambit of  groups consciously established for the purposes of literary discussion (for a history of the book club institution in the U.S. see Long \citeyear{long03}).  This is in keeping with the highly institutionalized perception of (most) literature (and reading as a cultural practice) as ``serious.''Conversations about such cultural fare as movies, however, serves to cement interaction in less formalized social settings \citep{fine77}.
 
\section{The stratification of culture-mediated social interaction}
\subsection{Estimation}
The key challenge in estimating the effects of socio-demographic markers on the propensity to engage in culture-based talk is tied to the peculiar nature of this dependent variable.  More specifically, the effect of any one socio-demographic characteristic on culture-based talk in any one domain is likely to be shaped by the fact that the distribution of individuals along a culture \emph{talk} continuum is simultaneously determined (in ways that we will attempt to disentangle below) by their distribution along a culture \emph{consumption} gradient (and vice-versa). Because it is likely that the same predictors of culture-talk will also figure as predictors of culture consumption, the issue of the mutual determination of talk and consumption has to be dealt in any attempt to model the socio-demographic heterogeneity in the propensity to engage in both activities.  

\subsection{Statistical model}
I specify the effects of social position on the propensity to engage in culture-based talk using a bivariate probit model \citep{ greene98}. This model is useful precisely when the analyst is faced with binary or categorically ordered dependent variables which may share a causal history or are simultaneously determined \citep{calhoun89}.  In the standard univariate probit specification for instance, the underlying propensity for each person to engage in culture-talk or aesthetic consumption would be conceptualized as an unobserved latent (normally distributed) variable ($y^{*}$).  The empirical realization ($y$) of the outcome variable is binary and equals one if the underlying latent propensity is above some threshold (e.g. $y^{*} \geq 0$) and zero otherwise. The probability of being above or below the threshold is stipulated to be a linear function of a vector of socio-demographic covariates ($\beta x$), and all the other factors that help us to predict the outcome not contained in this vector, are assumed to be independent and identically distributed \citep{long_freese06}.

\begin{center}
	Figure~\ref{fig3} about here
\end{center}

The path-model implied in the standard of the bivariate probit specification is shown in Figure~\ref{fig3}. The bivariate probit takes the same general form as the traditional univariate probit model, with the exception that the joint dependent variable is now considered to be the empirical realization of two latent underlying variables ($y^{*}_{1}$ and $y^{*}_{2}$) with random components presumed to follow a bivariate normal distribution.\footnote{This model is also referred to as the ``seemingly unrelated'' bivariate probit if the two equations contain different predictors, which is not the case in the analysis that follows.} The relationship between the latent normal variates ($y^{*}$) and the observed binary variables ($y$) is as follows:

\begin{equation}
y=g(y^{*})=\left\{
        \begin{array}{rl}
                0 &\mbox{if $y^{*} \leq 0$} \\ \\
                1 &\mbox{if $y^{*} > 0$}  \\
        \end{array} \right.
        \label{eq1}
\end{equation}

\begin{equation}
        \begin{array}{rl}
                y^{*}_{1i}=\beta_{1} x_{1i}+u_{i} \\  \\
                y^{*}_{2i}=\beta_{2} x_{2i}+v_{i} \\ \\
                corr(u,v)=\rho \\ \\
                u,v \approx N(0,1)  \\
        \end{array}
        \label{eq2}
\end{equation}

Where $\beta_{1}$\ and $\beta_{2}$ are vectors of unknown regression parameters to be estimated.  The explanatory variables in the model are presumed to be exogenous so that $cov(x_{1i},u_{i})=0$ and $cov(x_{2i},v_{i})=0$;  the errors are assumed to be independently and identically distributed as bivariate normal, with mean zero and standard deviation of one. The model allows for the correlation ($\rho$) between the exogenous errors for each equation ($u$ and $v$) (which is equivalent to the conditional tetrachoric correlation coefficient between the two outcomes) to be estimated from the data. In the current application, the $\rho$ parameter represents the \emph{shared} exogenous factors that predict both culture-talk and culture consumption and which are not included in the systematic component of the model. 

One problem that arises in the present case is that neither the culture-based talk and the culture-consumption indicators at hand are binary. Instead each is an ordered response variable with four categories: (1) Not at all, (2) Not very often, (3) Fairly Often and (4) Very Often. To accommodate this feature of the data into the bivariate probit modeling framework I rely on a further generalization of the bivariate probit model known as the bivariate ordered probit model \citep{calhoun89}.  

 In the bivariate ordered probit estimation procedure, each of the endogenous underlying latent variables ($y^{*}$) is now conceptualized as a latent normally distributed variable whose empirical realization is observed as bounded in ordered categories, following the constraints (in the four-category case):

\begin{equation}
y=g(y^{*})=\left\{\begin{array}{rl}
  1 &\mbox{if $y^{*}p \leq \tau_{1}$} \\
  2 &\mbox{if $\tau_{1} < y^{*} \leq \tau_{2}$}\\ 
  3 &\mbox{if $\tau_{2} < y^{*} \leq \tau_{3}$}\\ 
  4 &\mbox{if $\tau_{3} \leq y^{*}$} 
 \end{array} \right.
 \label{eq3}
\end{equation}

Where the unknown (to be estimated from the data) cutoffs satisfy the condition $\tau_{1}<\tau_{2}<\tau_{3}<\tau_{4}$; for purposes of normalization (and identification), $\tau_{0}$ is set to $-\infty$ and $\tau_{4}$ is set to $\infty$ \citep{long_freese06}.

\subsection{Variation in culture-based talk and culture consumption by social position}
Likelihood ratio $\chi^{2}$ statistics corresponding to the test of the hypothesis that a given block of variables is equal to zero, are shown in Table~\ref{tab1a}.  I specified a bivariate ordered probit model that included the same set of predictors for both the culture-talk and the culture-consumption equations. The predictors were chosen  using a combination of criteria geared towards substantive relevance and predictive power. I was of course also limited by the availability of socio-demographic indicators in the survey: for instance, the NSS contains no information on respondent's current occupation or family background (both of which are theoretically relevant dimensions of social differentiation). I ultimately decided to examine variation in culture-talk and culture-consumption across five broad axes of social differentiation: (1) Education (in Five categories), (2) Age (in seven categories), (3) Gender, (4) Income (in six categories)\footnote{Family income information was not available for about one fifth of the cases (20.8\%) across the two survey waves.  For these cases ($N=375$) I imputed their income by using the median family income category for that education/gender/marital status celll.  Including a dummy variable that equals one for the cases in which income was imputed from the data does not change the results.} (5) Race/Ethnic Identification (in Five Categories). The proportion of respondents in each of the categories of the predictor variables included in the model are shown in Table~\ref{tab3} of the Appendix. 

The coefficient estimates (and the associated t-statistic corresponding to the hypothesis that the regression weight corresponding to that variable is zero) for each social position indicator are shown in the left-most set of columns (for each cultural domain) in Tables~\ref{tab2a} and ~\ref{tab2b}.  A \emph{t}-value with an absolute value above 1.96, indicates that there is significant heterogeneity along that social dimension in the propensity to engage in culture-mediate interaction in comparison to the base category. Likelihood ratio $\chi^{2}$ statistics corresponding to the test of the hypothesis that a given block of variables (e.g. education, age) is equal to zero, are shown in Table 3a. The table also displays the share of the model $\chi^{2}$ that each block of variables is responsible for as a proportion of the sum of the $\chi^{2}$ accounted by all of the other socio-demographic variables in the model.  These percentages give a rough indication of the relative importance of that set of variables (in comparison to one another) in predicting the propensity of culture-based talk. The same is done for culture consumption in Table~\ref{tab1b}.\footnote{The bivariate ordered probit equations have all-categorical predictors in the right-side.  To make interpretation of the relevant effects more manageable, I omit the constant from the model and instead of using standard dummy coding, I constrain the coefficient estimates for each block of categorical variables to sum to zero.  This means that all of the categories (for those variables that have more than two categories) are included in the model.  Thus, there is no omitted reference category  (for a variable with $k$ categories, $\sum_{k}\beta{k}=0$). For variables with only two categories only one coefficient is shown, since the other one has the same absolute value but opposite sign.}



\begin{center}
	Table~\ref{tab1} about here
\end{center}

\begin{center}
	Figure~\ref{fig4} about here
\end{center}

As indicated by the the likelihood-ratio statistics in shown in Table~\ref{tab1a} and by the positive coefficient estimates corresponding to relevant differences across educational groups shown in Table~\ref{tab2a}, culture-based talk is unequally distributed by education across \emph{all} five domains.  Persons with more experience in the educational system are more likely to use culture in conversation than the less educated ($p<0.01$) \citep{collins75, dimaggio96}.   Naturally, education is not an equally powerful indicator of the likelihood to engage in culture-based talk across the different cultural realms.  Differentiation by education (in comparison to heterogeneity premised on age, gender, race and income) is preponderant in the case of culture-based talk involving literature and the arts, and is relatively less important in the case of movies, although it is still fairly substantial for Music (32\% of the total model $\chi^{2}$).  These results suggest that the highly educated are able to mobilize a relatively wide range of cultural resources in interaction, while those of low education are less likely to use culture (whether commercial or associated with the fine arts) as a topic in conversation. 

Looking at the coefficient estimates for the effects of education on culture consumption shown in Table~\ref{tab2b} I find that, consistent with almost all studies conducted on the subject of cultural participation, education-based differentiation in the likelihood of being an active culture consumer exists in all domains connected to the non-profit arts, and favors the educated (as frequent consumers) over the less educated (as frequent abstainers) \citep{baumol_bowen66, dimaggio_useem78a, bourdieu_darbel91}.  When it comes to relative rates of engagement in the plastic (55\% of $\chi^{2}$) and performing arts (35\%) and pleasure reading (45\%), education accounts for a substantial share of the explanatory power of the five axes of social differentiation considered here.  These results supports the proposition that education is one of the main predictors of frequency of engagement with the more artistically-oriented---restricted-production---domains, but not in the more ``commercial''---large-scale production--domains \citep{bourdieu84, bourdieu93a, lizardo_skiles09}.  The predictive power of education in determining who is an active culture-consumer is relatively muted (but still substantial) for Movies (24.3\% of $\chi^{2}$) and almost non-existent for Music  ($\chi^{2}=4.0, p=0.41$); this last result contrasts sharply  with the fact that there exist sharp differences in the use of music-culture in conversation for members of different educational groups.  

Differentiation linked to age categories in the underlying likelihood to use culture in interaction, where it exists, favors the young over the old:   younger persons are more likely to mobilize culture as a resource in conversation than those who are older. Age heterogeneity in the propensity to engage in culture-based talk is very salient for music ($\chi^{2}= 64.0, p<0.01. df=6$), but it is extremely salient in the case of movies ($\chi^{2}=221.1$). Differentiation by age in the propensity to use culture in interaction is in comparison very weak for the plastic ($p=0.07$), performing arts ($\chi^{2}=23.4, df=6$) and literature($\chi^{2}=15.6$), suggesting that the young are more likely than the old to mobilize cultural knowledge from the commercial arts, but are equally likely than their older counterparts to talk about the fine arts.  Overall we can conclude that the elderly are the group who are the least likely to use culture---in particular the commercial and popular arts---as a resource in conversation \citep{dimaggio87}, while younger adults are the most likely to do so especially when it comes to commercialized aesthetic domains \citep{frith98}.

Age-based differentiation in frequency of culture-consumption mirrors what we find in the case of culture-talk.  For Movies and Music, age is overwhelmingly the primary factor differentiating avid consumers from abstainers ($\chi^{2}=93.7, df=6$ for music and $\chi^{2}=211.8$ for movies), with the young being more likely to go out to the movies regularly or being a frequently music listener than older persons as is shown clearly by the coefficient estimates displayed in Table~\ref{tab2b}.  In contrast to the predictive ability of age in the realm of Movies and Music, there is no comparable differentiation across age-groups in terms of frequent engagement in the domains associated with literature and the non-profit arts, once we hold constant education, although a relatively moderate age-related differentiation effect can be ascertained for the performing arts ($\chi^{2}=35.7, p<0.01$). 

Gender differentiation in the underlying propensity to engage in culture-talk (favoring women over men) can be observed in all domains except for music ($\chi^{2}=2.7, p=0.10, df=1$), but the gender difference is relatively weak when it comes to the arts ($\chi^{2}=14.3, df = 1$) and is most salient in the domain of pleasure reading ($\chi^{2}=60.5$) \citep{long03, tepper00}, with gender accounting for about 34\% of the overall equation $\chi^{2}$ in this last case.   Gender-based differentiation in culture-talk is weaker than that premised on education and age for all domains, except for reading, where the effect of gender is comparable to the age effect and stronger than the income effect.  Differences between men and women in arts participation frequency are generally weaker than those found for culture-based talk.  Women are no more likely than men to engage the plastic arts ($p>0.05$) nor are they less likely to enjoy live drama ($p>0.25$). Women appear to be slightly more likely to listen to music frequently than men ($\chi^{2}=7.7, df=1$); they also appear marginally more likely to go to the movies frequently ($\chi^{2}=5.2$).  Consistent with the bulk of previous research on the matter, women are substantially more likely than men to read books for pleasure regularly ($\chi^{2}=37.6, p<0.01$) \citep{tepper00}; in this case, gender is the second most important differentiating factor (after education).

Within levels of education and income, race and ethnic differentiation in culture-talk propensities is generally weak, and where relatively strong effects can be found they appear to follow no substantively meaningful pattern. This suggests that access to culture as an interactional resource is primarily bounded by indicators of generalized status (such as education) and life-course stage and not by categorical membership in certain status groups premised on race (even as it is very likely  that these cultural discussion networks to be segregated by race \citep{carter03, dimaggio_ostrower90}).  The same goes for racial and ethnic differences in culture consumption which also tend to be small and weak.  

I find that culture-based talk differentiation premised on family income is generally weak.   Where income-based differences in the likelihood of engaging in culture-talk exist they show a pattern similar to stratification based on education: low-income groups are less likely to engage in culture-based talk than economically advantaged respondents, as shown in Table~\ref{tab2a}.   Income-based differences in culture-talk propensities are strongest in the case of Drama ($\chi^{2}=23.5, df=5$), but they are also statistically discernible for Music and Movies as well, even if their substantive significance is low. Differences across high and low income groups are small (in terms of magnitude) for Music, Movies, Literature and the Arts.   The pattern of comparatively weak income effects is consistent with the view that age, education and gender provide more powerful axes of status based stratification (based on access to informal networks of acquaintanceship) than does raw differences in income streams  \citep{dimaggio_useem78a}.   Income is generally a better predictor of consumption frequency than it is of culture-based talk.  Income-based differences in participation display the same pattern as that evident in culture-based talk: members of high-income groups participate in the arts at higher rates than persons classified as low income.  Income differentiation in cultural participation is null for Music and Literature but it is somewhat important for the two domains associated with activities likely to involving evenings outside of the house: Movies ($\chi^{2}=51.7, p<0.01$) and Theater ($\chi^{2}=36.3, p<0.01$).  

\subsubsection{Recursive bivariate probit estimation}
The bivariate probit specification model presumes that net of the shared correlation among the exogenous shocks ($\rho_{21}$), culture-consumption  and culture-talk are mutually exogenous.  However, as shown in Figure~\ref{fig2}, this assumption is in all probability too strong:  frequent consumers in a given cultural-domain should be expected to use that particular type of cultural content in interaction (even after adjusting for the relevant socio-demographic covariates).   In terms of the current modeling framework, this means that culture consumption has an \emph{endogenous effect }on the frequency of culture-talk, implying that the relevant causal system is closer to that depicted in Figure~\ref{fig4}.   Looking at the figure, it is evident that if ($\delta_{21}>0$) then the respondent's socio-demographic characteristics, in addition to having direct effects on culture-talk (through ($\beta_{1}$), also have \emph{indirect effects} (as mediated through the $\delta_{21}\beta_{2}$ pathway) on the same outcome.  

The question becomes:  what share of the effects of age, education, gender and income shown in Table~\ref{tab2a}, are due to the fact that persons of a different education, age or gender are more likely to be frequent participants in a given domain?  For instance, it is possible that \emph{all} of the culture-talk differences across education or age levels are due to the existing consumption heterogeneity across age and  education groups.  This is the case in which there are no direct effects of social position on culture-talk, since they are entirely mediated via culture consumption (e.g. $\beta_{1}^{educ}=0|\delta_{21}, \beta_{2}^{educ}$).  Relatedly, if the bulk of the direct effects of socio-demographic markers on culture-talk remain even after controlling for the endogenous effect of consumption, this is an indication that other axes of social differentiation---not linked to the probability of being a frequent consumer of the arts---are implicated in generating cross-group differences in culture-talk propensities.

In order to shed light on these issues, I specify a recursive bivariate probit model \citep{greene98} with an endogenous effect of consumption on culture-talk.  The recursive bivariate ordered probit with an endogenous regressor can be written as:

\begin{equation}
        \begin{array}{r r r}
		y^{*}_{1i} &  = & \delta_{21}y_{2i} +\beta_{1} x_{1i} + u_{i} \\ 
		y^{*}_{2i} &  = & \beta_{2} x_{2i} + v_{i} \\
        \end{array}
        \label{eq4}
\end{equation}

With everything else as in Equation~\ref{eq2}, and with the implicit restriction: $\delta_{12}y_{1i}=0$ \citep{simon54}.  This model specifies a direct effect of one of the endogenous variables on the other one, with consumption frequency used as a predictor of the rates of culture-talk ($y_{2} \rightarrow y_{1}$).  This makes this model a member of the class of ``recursive'' bivariate probit models \citep{monfardini_radice08}.   These models are especially useful when dealing with the problem of estimating directional effects with simultaneously determined endogenous outcomes \citep{greene98}.  What most analysts usually miss is that in these models there is no need to include additional ``instrumental variables'' or to rely on other so-called ``exclusion restrictions'' to estimate the endogenous effect (this means that the $x_{1}$ and $x_{2}$ vectors can contain the same set of predictors) \citep{monfardini_radice08}.  Instead, one can simply ``plug in'' one of the jointly determined variables as a right-hand side predictor of the other---as depicted graphically in Figure 4 \citep{greene98}.\footnote{As long as it is understood that  $y_{2i}$ enters the model as \emph{observed} and not it its latent, continous form.} The analyst can then use the usual Wald and t-tests to determine whether a model that specifies an effect of $y_{2}$ on $y_{1}$ gives a better representation of the data-generating process than one in which the outcomes are presumed to be mutually exogenous.\footnote{Including exogenous predictors in the interaction equation that are excluded from the taste-equation, while not strictly necessary for identification, may help make inference more robust to distributional misspecification \citep{monfardini_radice08}.}

\subsubsection{Results}
Coefficient estimates of the endogenous effect of consumption on culture-talk obtained from the recursive bivariate probit estimation are shown in Table~\ref{tab3}.\footnote{In the model that I estimated, when the ordered categorical variable corresponding to consumption frequency ($y^{*}_{2i}$) appears on the right hand side (as a predictor of talk). I constrain the endogenous consumption effect to be linear by linear \citep[370]{agresti02}; in addition, if the coefficient estimate for the endogenous effect is found to be larger than zero, then I constrain the $\rho$ estimate to be zero.}  The results demonstrate that, as expected, consumption has a positive endogenous effect on culture-talk in all five domains, even after holding adjusting for compositional differences due to socio-demographic location.  The endogenous effect of consumption is strongest for the Arts ($\delta_{21}^{Arts}=0.76, t=20.78$) and is weakest in the case of music ($\delta_{21}^{Music}=0.34, t=10.67$).   

To examine how the effect of socio-demographics characteristics on culture-talk change after accounting for the endogenous effects of consumption, I compute the \emph{marginal effect} corresponding to each binary indicator of social position \citep{[126]long_freese06}.  The marginal effect gives the average shift in the probability of observing a respondent with a given level of culture-talk given that we ``turn-on'' a particular socio-demographic factor from its null value of zero to one.  For instance, switching on the ``post-graduate degree'' indicator, results in a positive change in the probability of engaging in culture talk very often; switching on the ``65 or older'' indicator results in a negative change in the probability of talking to a friend or acquaintance about the movies very often.  I calculate marginal effects for each socio-demographic variable both for the standard bivariate probit model and for bivariate probit with the endogenous consumption effect.  The absolute size difference in the marginal effects gives us an estimate of how much of the effect for that socio-demographic variable is mediated through consumption.  

A bar graph for the marginal effect of each socio-demographic variable in the bivariate and the endogenous model for all five domains is shown as Figure~\ref{fig5}.  The marginal effect for each variable in the bivariate model is shown in white, while the marginal effect in the corresponding model with the endogenous consumption effect is shown in black.  The difference in the height of the bars is proportional to the extent to which the effect of that variable is adjusted when we account for the indirect of socio-demographic factors through of consumption.\footnote{The marginal effects shown in the graph correspond to the average shift in the probability of the respondent falling into either the ``fairly often'' and ``very often'' categories for each cultural domain (this is obtained by adding together the estimates of the marginal effects of falling into third and fourth categories of the culture-talk frequency response).}

Considering the change in magnitude of the socio-demographic effects, I find that most cross-group differences are\emph{ partially}, but not \emph{entirely}, mediated by socio-demographic heterogeneity in consumption frequency.  For the most part specifying a direct effect of consumption on culture-talk drives the absolute magnitude of the coefficient estimates for the majority of the socio-demographic covariates downwards, which is consistent with a partial mediation account.  However, even after adjusting for the fact that the educated are more likely to engage the Arts and literature, I find that education is still a substantively significant predictor of culture-talk in all five domains, and that age is an important marker of who engages in culture-talk involving music and movies (but not literature and the arts). In a similar way, gender continues to be an important predictor of who engages in culture-talk involving literature, and gender differences in other realms are essentially unaffected by the inclusion of the endogenous effect of consumption.   

Naturally, the biggest deflation in the marginal effect of belonging to the highest education group (postgraduate degree) can be appreciated in the most socio-demographically exclusive consumption domains: Literature and the Arts; in the case of Music there is no much difference between the bivariate and the endogenous model in the marginal effects of education  This is to be expected, since there is no education-linked heterogeneity in frequency of music listening (and thus no indirect effect of education through consumption). In the case of age, we can see that part of the age-related differences when it comes to culture-talk involving movies are attributable to the fact that older persons go to the Movies less frequently than the young.  However, even after holding constant consumption, age continues to be a substantively significant predictor of culture-talk for film and music.  

\subsection{Culture-based talk across taste cultures}
To answer the third research question that I posed above, I use information on the respondent's favorite musical genre-cluster to ascertain the effect of membership in distinct taste cultures on the propensity to engage in culture-based talk.  I use the respondent's self-report on aesthetic preferences as an empirical indication of the taste-culture that the respondent belongs to.\footnote{Most respondents have only one favorite. Only a small minority chooses more than two.}  The ``intellectualization'' hypothesis suggests that it is those persons who display an ``aesthetic disposition'' \citep[3]{bourdieu84} toward those genres most clearly institutionalized as ``art'' that will be more likely to engage in culture-based talk.   Music emerges as an appropriate case in this respect since, as we we have already seen, members of all educational groups are about equally likely to report listening to music very often; yet respondents with a bachelors or a post-graduate degree are substantially more likely to report \emph{talking} about music with their friends while their low education counterparts are less likely than average to do so. 

This pattern of results can be accounted for if there exists structured variation in culture-based talk in the music domain across persons who belong to different musical \emph{taste cultures} \citep{gans99}.  Respondents more likely to display a preference for those genre categories more likely to attract high-status individuals are also more likely to generate culture-based talk than those who display a preference for genres more likely to attract individuals of lower educational status. In this manner, heterogeneity among music listeners in terms of allegiance to distinct (but hierarchically ordered) taste cultures may partially account for why we observe educational heterogeneity in terms of the deployment of musical culture as a resource in interaction.  This heterogeneity can exist without a corresponding form of education-based stratification in terms of the frequency of listening to music.   

The statistical model that I use to test the intellectualization hypothesis is the same as the one I specified to estimate the endogenous effect of culture-consumption on culture-talk.  I specify a recursive bivariate ordered probit model with correlated errors, with frequency of culture-talk ($y_{1}=$1{\ldots}4) and favorite musical genre-cluster ($y_{2j}=1$ if the respondent reports liking genre $j$ best and $y_{2j}=0$  otherwise) as the joint dependent variable.  The model includes an endogenous effect ($\delta_{21}$) of music-taste culture membership on the frequency of engaging in culture-talk involving music (see Figure~\ref{fig4}).\footnote{This statistical model differs from the one used to investigate the effects of consumption on culture-talk in one main respect. Instead of both variables being ordered categorical outcomes (resulting in the recursive bivariate ordered probit model) one of the dependent variables (favorite musical genres) is binary ($y_{2}=1$ if the respondent likes that music and $y_{2}=0$ otherwise) and the other one (frequency of culture-based talk) has multiple ordered categories. This makes this model a recursive bivariate ``semi-ordered'' probit model \citep{armstrong_mcvicar02}.}  We should expect education (and other socio-demographic covariates) to have \emph{indirect} effects on culture-talk frequency through the cultural taste pathway, comparing the estimates of the education effect in the recursive model from those obtained from the simple bivariate probit model gives us a sense of how much of the direct effect of education is due to educational heterogeneity in taste culture membership.   

\subsubsection{Results}
Selected coefficient estimates from the recursive bivariate semi-ordered probit model are shown in Table~\ref{tab4}.\footnote{To estimate the models, I proceeded as follows:  I first estimated the model without the recursive effect.  If a Wald test showed that the estimate of $\rho$ was significantly different for zero (indicating that taste and music-talk were jointly determined and thus not exogenous), then I re-estimated the recursive model with the $\rho$ parameter allowed to vary freely.  If the Wald test indicated that $\rho$ was statistically indistinguishable from zero, then I estimated the recursive model with this parameter constrained to be zero. See \citet{greene98} for a similar analytic strategy.} Although the fact that reports on taste are forced to conform to the highly aggregate (and somewhat internally heterogeneous) genre categories shown in the table is not ideal, the overall pattern of results conforms to expectations.\footnote{Unfortunately, the wording for the items eliciting the respondent's favorite group of musical genres is different in the two waves of the NSS. While the 2004 item comes closer to an actual elicitation  of the respondent's favorite music (\emph{Do you have any favorite types of music, for example, rock and roll, popular music, country music, religious music, jazz, or opera?}), the 2002 item's wording makes it closer to a probe for which musical genres are closer to the respondent's sense of cultural (or personal) identity (\emph{What kind of music is that [that you think of as very important to you, or to who you are]?}).  Because the two forms of the question are certainly not  comparable, I restrict myself to the 2004 data ($N=758$).}  

As shown in Table~\ref{tab4a}, members of the Classical music or Opera taste culture display a substantially greater likelihood of engaging in culture-based talk involving music ($\beta=0.83, t=13.77$), net of all other socio-demographic predictors.  Those who report liking Jazz or Blues are more likely to use music as a resource in interaction  ($t=3.78$), as are members of taste cultures expressing a liking for Rock and Pop (although the magnitude of the effect is markedly lower: $\beta=0.11, t=2.55$).   For members of taste cultures featuring lower-status genres, the taste-culture effect approaches zero (in the case of Punk and Heavy Metal, however, the coefficient estimate of $\delta$ is positive and the magnitude is similar to that for Rock and Pop although it does not reach statistical significance).  Members of taste cultures expressing a preference for Christian and Gospel are no different from the average respondent, while members of taste cultures that report Country and Folk genres as their favorite music actually are \emph{less likely} than other respondents to engage in culture-based talk ($\beta=-0.67, t=-8.64$).   

In addition, note that---as shown in the lower panel of table~\ref{tab4a}---after including the endogenous effect of membership in the most elite taste culture, the direct effect of education on culture-talk is greatly muted, to the point that it is no longer statistically significant at conventional levels ($|t|<1.96$).  This is in stark contrast to the model including the endogenous effect of consumption in which the direct effects of education remained largely intact (as shown in Figure~\ref{fig5}).  This suggests that the educational differentiation in culture-talk involving music is entirely mediated by the fact that the higher-educated are more likely to associated themselves with higher-status taste cultures \citep{gans99, bourdieu84}. 

These results are consistent with the hypothesis that affiliation with taste cultures associated with those musical genres most likely to have already institutionalized organizational and critical resources for intellectualization increases the likelihood that members of those taste culture will engage in culture-based talk. Respondents who express a preference for genres which are less likely to feature institutionalized resources for intellectualization are relatively less likely to mobilize musical culture in interaction\footnote{This statement should not be taken to imply that music plays \emph{no role} in the person's life of groups marked by low educational status, as recruitment for the generation of sociability is only one of many ways in which music can structure everyday lived experience \citep{denora00}.}   These results also help us to partly explain why we observe stratification of culture-based talk in the music domain in the absence of education-linked differences in terms of consumption frequency.  

Table~\ref{tab4b} shows the coefficient estimates (and \emph{t}-statistics) corresponding to the estimates of education-based heterogeneity on the probability of belonging to each taste culture.  The table shows that the likelihood that membership in a taste culture is associated with the production of higher rates of culture-talk (Table~\ref{tab4a}) is positively associated with: (a) the relative likelihood that members of that taste culture hail from the ranks of the highly educated and (b) the likelihood that that taste culture repels members of low education groups (in the case of Classical, Opera and Jazz).  The predictive strength of taste-culture membership in culture-talk propensity also follows the direction of differentiation by educational group on taste, with those taste-cultures weakly differentiated by education (e.g. Rock and Pop) having relatively weak effects on culture-talk probabilities, and those taste-cultures strongly differentiated by education (with Opera/Classical versus Folk/Country at different ends of the spectrum) having strong (but naturally divergent) effects.\footnote{The only exception to the pattern is the Punk/Heavy Metal subculture, which recruits from low education groups (and repel high education respondents), without resulting in a negative effect on culture-based talk.}

\section{Discussion and Conclusion}
\subsection{Contributions}
\subsubsection{The theoretical importance of culture-talk}
This paper contributes to the study of culture and stratification in several ways.  Analytically, I highlight the theoretical and empirical importance of the phenomenon that I have referred to as \emph{culture-based talk}. I began with the Simmelian legacy on the importance of seemingly trivial conversational content for the emergence of ``sociability'' (or the ``play form'' of association) and connected it to the more contemporary stress by cultural stratification theorists on ``entertainment talk'' as the cement that holds status groups together as micro-relational communities of acquaintance and informal membership \citep{collins75, dimaggio87, purcell07, lizardo06a}.  I have argued that the deployment of cultural knowledge in concrete interaction is an important mechanism linking culture to status-based stratification.  In spite of its centrality, it has been a relatively under-studied phenomenon. This paper contributes to the study of culture and stratification by examining the socio-demographic distribution and cultural correlates of the joint tendency to talk to friends about culture consumption experiences.  

\subsubsection{Mapping the social ecology of cultural talk}
Empirically, I advance research on the role of culture-based talk as a constitutive component of lifestyle-based stratification in contemporary post-traditional societies.  I do this by examining the inter-relationship between aesthetic consumption, culture-talk and cultural taste. The results show that culture-based talk is a relatively common experience for the large majority of Americans, with almost four fifths reporting engaging in some form of culture-based talk using cultural knowledge from the five domains included in the NSS. The results also show however, (consistent with prior theory in cultural stratification) that culture-based talk is stratified by the same markers of social position (Age, Gender) and social status (Education) that previous research has shown also characterize audience segmentation in the corresponding cultural domains.  

In particular, I show that culture-talk in those domains characterized by large-scale commercial culture (Movies, Music) is strongly stratified by education (with the highly educated generating more culture-based talk) and in particular by age (with the young being more likely to use such culture in conversation). Culture-talk recruiting cultural knowledge from the higher-status ``non-profit'' arts is largely stratified by educational attainment, but only weakly differentiated across age-groups (net of education).  Literature is distinctive in that in addition to exhibiting the expected substantial rates of stratification by education it is also strongly stratified by gender (although all domains except for music show a female advantage in culture-based talk) as suggested by recent research on compositional differences of the ``reading class.''  

These results provide baseline expectations for the likely distribution of cultural content across socio-demographic groups.  For instance inter-generational segmentation in informal sociability is likely to be informally reproduced and generated by age-based segmentation on what type of culture is deployed in interaction \citep{milner04}, a mechanism that also serves to maintain divisions across educational groups \citep{collins75}.   By showing that culture-based talk is relatively commonplace but also socially structured in predictable ways, this paper takes a first step towards the empirical specification of an important but so far under-analyzed phenomenon in the sociology of taste.  

\subsubsection{The interactional basis of the omnivore phenomenon}
This paper also contributes to our understanding of the phenomenon of cultural omnivorousness as a primary axis of cultural stratification \citep{peterson92, peterson_kern96}.  In particular the results reported above help to specify how the omnivore disposition manifests itself as an \emph{interactional} phenomenon \citep{ollivier08}. Three sets of findings are particularly pertinent in this respect.  First, I show that relative rates of culture-talk increase with education in all domains, whether commercial, literary or artistic. This converges nicely with the finding that education is also the one of the best socio-demographic predictors of cultural omnivorousness \citep{warde_etal08, peterson_rossman08, lopez-sintas_katz-gerro05}.   Taken together, these findings suggest that cultural omnivores are not only more likely to consume a wide variety of symbolic goods, they are also more likely to mobilize the cultural knowledge so-gained in interaction \citep{peterson92, dimaggio96, lizardo06a}. 

Second, the results indicate that the educational stratum most likely to recruit cultural knowledge in interaction consists of those who have been bestowed a postgraduate degree. This is consistent with the contention that culture-talk is an interactional resource characteristic of high-status groups in contemporary American society \citep{collins04}. This is also consistent with prior research that shows that high aesthetic-consumption groups, are also those more likely to possess richer networks of strong ties (which provide closure).  Such markers of social position as income, race and ethnicity and gender in contrast, are distinctive for the fact that they do not serve as powerful differentiating factors in predicting the propensity of culture-based talk. This is consistent with the view that the educational system is the primary institution responsible for providing culture-consuming publics with the dispositions to treat cultural goods as things that can be classified, sorted and talked about \citep{bourdieu93b, dimaggio00, lizardo08}.  

Finally I find that, for the most part, the stratification of culture-based talk in each domain mirrors the stratification of aesthetic consumption in that same consumption realm. Nevertheless, important discrepancies are evident.   In particular, the realms more deeply connected to ``commercial'' culture (movies and music) exhibit much weaker stratification by education in terms of consumption than they do in terms of talk.  High-status individuals are able to generate culture-based talk using knowledge procured from so-called ``popular'' cultural domains even as their consumption rates in these domains do not differ as dramatically from the less educated as they do in the fine arts realm (literature and the non-profit arts). This is another way in which the cultural advantage associated with omnivore taste, is made concrete as the facility of higher status persons to mobilize cultural knowledge taken from heterogeneous domains in interaction \citep{ollivier08}, thus serving as bridges across ``cultural holes'' \citep{lizardo06a, pachuki_breiger10}.

\subsubsection{Heterogeneity in rates of culture-talk across taste groups}
The results reported above show that a model that incorporates recent research and theory on ``intellectualization'' in cultural domains institutionalized as artworlds, generates a testable prediction: those persons more likely to be oriented toward high-status culture that will be more likely to generate culture-based talk. The results show that persons exhibiting high-status taste patterns in the domain of music, are more likely to generate culture-based talk within that domain.  This explains why we observe education-based stratification in the likelihood of generating culture-based talk in this domain in the absence of education-linked differences in the frequency of listening to music \citep{denora00}. These results contribute to our understanding of how the ``aesthetic disposition'' is mobilized in everyday interaction as a higher likelihood to engage in culture-based talk across domains. 

An important question for future research is whether a similar pattern can be observed in other domains associated with the commercial arts (e.g. Movies) or literature. The intellectualization hypothesis predicts that it is precisely those audiences more likely to gravitate towards the most ``intellectualized'' sub-genres (e.g.  in the case of movies, such genres as foreign, ``independent'' or documentary film) that are also more likely to use culture as an interactional resource.    Nevertheless, we must leave open the possibility that future studies that include a wider range of cultural practices and domains may lead to conclusions that heavily qualify the proposition connecting intellectualization and genre-legitimacy to culture-talk.  For instance within the realm of television-programming, it may be those genres endowed with the least amount of critical legitimacy (e.g. reality television) that generate the most culture-based talk, especially among the young.  The very limited research done on ``sports talk'' \citep{purcell07} suggests however, that ``intellectualization'' in the realm of sports (e.g. produced by adopting the critical vocabulary produced by sports journalists) does facilitate sports-based talk among educated sports enthusiasts, as suggested by the above analysis. 
 
\subsubsection{Rethinking Important Matters}
The pervasiveness and social distribution of the phenomenon of culture-based talk has important implications for the ways in which data about social connectivity are collected and conceptualized.  The two large-scale efforts to collect data about patterns of network connectivity in the American population have relied on the ``important matters'' name generator as their main methodological tool \citep{burt84, marsden87, mcpherson_etal06}. The presumption here is that there is some sort of homology between the alleged \emph{importance of the content of conversation and the importance of the tie in question}.  Importance here being conceptualized both as 1) theoretical importance for the social scientist,--e.g. the status of the social ties identified in this manner as an measure of the extent to which the individual is ``connected'' \citep{burt90}; and 2) practical importance for the person; or the extent to which the connection identified in this manner give us access to a meaningful subset of the relations that a person would draw on in time of emotional or financial help \citep{wellman_wortley90}.

But given the social distribution of culture -mediated interaction---and its intuitive connection with certain socio-demographic markers---it is perilous to suppose that the ``importance'' of the ``matters'' discussed has any connection to the ``importance'' of the contact with which they are discussed.  In fact, given the theoretical framework sketched here, we should expect little or no connection between the objective importance of the network content and the object (or social-scientific) importance of the ties connected to to those contents.  If Simmelian ``pure sociability'' is more likely to be enacted among those contacts who provide the most \emph{intrinsic} satisfaction we should expect the connection to even be negative (e.g. meaningful and subjectively important ties are sustained through seemingly ``trivial'' and ``unimportant'' cultural content).  This becomes more important when we realize that, as Simmel would have predicted, conversation involving ``important matters'' is too sporadic and inconsistent a phenomenon with which to judge whether persons are connected or not \citep{bearman_parigi04}.

A substantial number of people (including a larger proportion of married persons) report simply not having anybody with which to discuss ``important'' matters with \citep{bearman_parigi04}). Should we conclude therefore that most persons are ``disconnected'' and experience little sociability \citep{mcpherson_etal06, putnam00}?  I would suggest that a more theoretically sophisticated understanding of what ``sociability'' is all about suggests otherwise. ``Important matters'' name generators may not be getting at either ``important'' content  or important ties (if the social-scientific interest is to uncover those ties that are meaningful for the individual).  This becomes even more of an issue when we consider that interaction based on important matters, may be more steeply stratified by markers of social status, such as education, than culture-based talk \citep{bourdieu84, mcpherson_etal06, marsden87, fischer82}.

In this respect, the ``important matters'' name-generator is a tough ``test'', that only a few privileged persons endowed with certain forms of linguistic (or political) capital may pass \citep{bourdieu84}. Important matters name generators confound stratification based on aptitudes orthogonal to social competence and ability to form connections with stratification criteria that directly impinge on the persons' ability to in fact create and sustain social relationships \citep{lin99, dimaggio96, lizardo06a}.  If the framework outlined here is on the right track, \emph{cultural matters may matter more than important matters} as the content with which to sustain (and thus measure) meaningful ties.  Thus, querying persons as to whom they discuss the ``trivial'' matters related to the cultural content that makes up entertainment talk (e.g. sports, fashion, cuisine, television, movies, the arts), maybe a more sensible strategy to generate data on personal social connectivity, and to  draw conclusions about the fate of this connectivity in the contemporary situation.

\singlespacing
\newpage
\bibliographystyle{asr}
\bibliography{N:/Private/zotero-library}
%\bibliography{C:/zotero-library}

%Figures%
\newpage
%%\begin{landscape}
\begin{figure}[t]
\centering
\includegraphics[width=8.5in, height=5.5in]{Univariate-Interaction.png}
\caption{Univariate distribution of culture-based talk probabilities (whole sample), 2002 and 2004 National Social Survey, Arts and Culture Supplement.}
\label{fig1}
\end{figure}  
%%\end{landscape}

\newpage
%%\begin{landscape}
\begin{figure}[t]
\centering
\includegraphics[width=8.5in, height=5.5in]{Univariate-Interaction-High-con.png}
\caption{Univariate distribution of culture-based talk probabilities (culturally active subsample), 2002 and 2004 National Social Survey, Arts and Culture Supplement.}
\label{fig2}
\end{figure}  
%%\end{landscape}

\newpage
\begin{figure}[t]
\centering
\includegraphics[width=6.5in]{sur.png}
\caption{Path model corresponding to the Bivariate Ordered Probit Model, where $y_{1}$ (culture-based talk) and $y_{2}$ (culture consumption) are simultaneously determined endogenous variables}
\label{fig3}
\end{figure} 

\newpage
\begin{figure}[t]
\centering
\includegraphics[width=6.5in]{recursive.png}
\caption{Path model corresponding to the Recursive Bivariate Semi-Ordered Probit Model, where the first endogenous variable $y_{1}$ (frequency of culture-based talk involving music) is an ordered categorical variable and the second endogenous variable $y_{2}$ (favorite musical genres) is a binary variable. In the model, $y_{2}$ helps to predict $y_{1}$, via the $\delta_{12}$ pathway.}
\label{fig4}
\end{figure} 

\newpage
%%\begin{landscape}
\begin{figure}[t]
	\includegraphics[width=9in]{marg-effects.png}
\caption{Marginal effects estimates of the probability of engaging in culture-talk either ``fairly often'' or ``very often'' for five cultural domains, 2002 and 2004 National Social Survey}
\label{fig5}
\end{figure}
%%\end{landscape}
     
%Tables%
\newpage
%%\begin{landscape}
\begin{table}[t]
\caption{Wald $\chi^{2}$ statistics and \emph{p}-values corresponding to the null-hypothesis that a given set of coefficients is jointly equal to zero, for the consumption equation of bivariate ordered probit models across five aesthetic domains, 2002 and 2004 National Social Survey.}
\subfloat [] {
\begin{tabular*}{\hsize}{@{\extracolsep{\fill}} l*{15}{c}} 
\hline 
\hline 
\multicolumn{15}{c}{culture-talk}\\
&\multicolumn{3}{c}{Educ. (\emph{df}=5)}&\multicolumn{3}{c}{Age (\emph{df}=6)}&\multicolumn{3}{c}{Gender (\emph{df}=1)}&\multicolumn{3}{c}{Race, Eth. (\emph{df}=4)}&\multicolumn{3}{c}{Income (\emph{df}=5)}\\
&\multicolumn{3}{c}{$H_{0}: \beta_{k}=0$}&\multicolumn{3}{c}{$H_{0}: \beta_{k}=0$}&\multicolumn{3}{c}{$H_{0}: \beta_{k}=0$}&\multicolumn{3}{c}{$H_{0}: \beta_{k}=0$}&\multicolumn{3}{c}{$H_{0}: \beta_{k}=0$}\\
& $\chi^{2}$  & \emph{p} & \% & $\chi^{2}$ & \emph{p}& \% &  $\chi^{2}$ &  \emph{p}& \% &  $\chi^{2}$ &\emph{p}& \% &  $\chi^{2}$ &  \emph{p}& \% \\ 
\hline   
Movies&31.3&0.000&0.102&221.1&0.000&0.718&27.4&0.000&0.089&0.3&0.990&0.001&12.0&0.034&0.039\\
Music&51.3&0.000&0.323&64.0&0.000&0.403&2.7&0.101&0.017&7.7&0.103&0.049&11.7&0.039&0.074\\
Books&66.3&0.000&0.369&15.6&0.016&0.087&60.5&0.000&0.337&9.0&0.060&0.050&10.1&0.071&0.056\\
Arts&78.0&0.000&0.540&11.7&0.068&0.081&14.3&0.000&0.099&8.6&0.073&0.059&8.6&0.128&0.059\\
Plays&62.5&0.000&0.353&23.4&0.001&0.132&24.5&0.000&0.138&12.5&0.014&0.071&23.5&0.000&0.133\\
\hline 
\end{tabular*}
\label{tab1a}
} \\

\subfloat [] {
\begin{tabular*}{\hsize}{@{\extracolsep{\fill}} l*{15}{c}} 
\hline
\hline
\multicolumn{15}{c}{Culture Consumption}\\
&\multicolumn{3}{c}{Educ. (\emph{df}=5)}&\multicolumn{3}{c}{Age (\emph{df}=6)}&\multicolumn{3}{c}{Gender (\emph{df}=1)}&\multicolumn{3}{c}{Race, Eth. (\emph{df}=4)}&\multicolumn{3}{c}{Income (\emph{df}=5)}\\
&\multicolumn{3}{c}{$H_{0}: \beta_{k}=0$}&\multicolumn{3}{c}{$H_{0}: \beta_{k}=0$}&\multicolumn{3}{c}{$H_{0}: \beta_{k}=0$}&\multicolumn{3}{c}{$H_{0}: \beta_{k}=0$}&\multicolumn{3}{c}{$H_{0}: \beta_{k}=0$}\\
& $\chi^{2}$  & \emph{p} & \% & $\chi^{2}$ & \emph{p}& \% &  $\chi^{2}$ &  \emph{p}& \% &  $\chi^{2}$ &\emph{p}& \% &  $\chi^{2}$ &  \emph{p}& \% \\ 
\hline    
Movies&24.3&0.000&0.073&211.8&0.000&0.631&5.9&0.015&0.018&15.5&0.004&0.046&51.7&0.000&0.154\\
Music&4.0&0.408&0.032&93.7&0.000&0.747&6.2&0.013&0.049&12.1&0.017&0.096&5.0&0.411&0.040\\
Books&59.9&0.000&0.448&14.3&0.026&0.107&37.4&0.000&0.280&11.9&0.018&0.089&3.2&0.673&0.024\\
Arts&100.4&0.000&0.552&7.3&0.293&0.040&2.5&0.110&0.014&1.1&0.887&0.006&19.3&0.002&0.106\\
Plays&60.1&0.000&0.345&35.7&0.000&0.205&0.3&0.614&0.001&5.5&0.244&0.031&36.3&0.000&0.208\\
\hline 
\end{tabular*}
\label{tab1b}
}
\label{tab1}
\end{table}
%%\end{landscape}

%%\begin{landscape}
\begin{table}[t]
\caption{Coefficient estimates of the effects of socio-demographic covariates on rates of culture-talk obtained from a bivariate probit model with frequency of culture-talk and frequency of culture consumption within five domains as the joint as the joint outcome, 2002 and 2004 National Social Survey}
\begin{tabular*}{\hsize}{@{\extracolsep{\fill}} l*{5}{c}}
\hline
\hline
&\multicolumn{5}{c}{Talk to someone you know about:}\\
&Movies&Music&Books&Arts&Performances\\
\hline
Education (Less than H.S.=1)  &-0.34 (-4.08)&-0.34 (-4.11)&-0.37 (-4.46)&-0.22 (-2.48)&-0.24 (-2.82)\\
Education (High School=1)     &-0.11 (-2.30)&-0.22 (-4.48)&-0.27 (-5.34)&-0.37 (-6.89)&-0.31 (-6.08)\\
Education (Some Coll.=1)      &0.06 (1.14)&0.11 (2.14)&0.12 (2.44)&0.02 (0.34)&0.05 (0.99)\\
Education (College=1)         &0.09 (1.69)&0.15 (2.79)&0.22 (4.03)&0.09 (1.68)&0.12 (2.18)\\
Education (Postgraduate=1)    &0.30 (4.38)&0.31 (4.51)&0.30 (4.29)&0.47 (6.75)&0.39 (5.58)\\
Age (18-24=1)                 &0.53 (7.01)&0.44 (5.91)&0.20 (2.77)&0.11 (1.37)&0.22 (3.00)\\
Age (25-29=1)                 &0.37 (4.48)&0.20 (2.50)&0.08 (1.02)&0.09 (1.09)&0.10 (1.27)\\
Age (30-34=1)                 &0.06 (0.76)&-0.08 (-0.97)&0.01 (0.17)&0.01 (0.09)&0.01 (0.13)\\
Age (35-39=1)                 &0.16 (2.01)&-0.01 (-0.12)&-0.07 (-0.87)&0.09 (1.17)&-0.01 (-0.12)\\
Age (40-49=1)                 &-0.00 (-0.02)&-0.10 (-1.65)&-0.03 (-0.43)&-0.04 (-0.67)&0.02 (0.35)\\
Age (50-64=1)                 &-0.31 (-5.78)&-0.14 (-2.61)&-0.03 (-0.58)&-0.08 (-1.35)&-0.15 (-2.67)\\
Age (65+=1)                   &-0.81 (-12.84)&-0.32 (-5.13)&-0.18 (-2.87)&-0.18 (-2.78)&-0.20 (-3.14)\\
Fam. Income ($<20K=1$)          &-0.12 (-1.69)&-0.18 (-2.58)&-0.11 (-1.56)&-0.07 (-0.90)&-0.27 (-3.75)\\
Fam. Income (20K to 35K=1)   &-0.14 (-2.40)&-0.04 (-0.73)&-0.09 (-1.50)&-0.06 (-1.08)&-0.08 (-1.37)\\
Fam. Income (35K to 50K=1)   &-0.06 (-1.18)&-0.06 (-1.13)&-0.03 (-0.50)&-0.08 (-1.51)&-0.10 (-1.91)\\
Fam. Income (50K to 75K=1)   &0.07 (1.28)&0.08 (1.46)&0.13 (2.30)&0.07 (1.25)&0.06 (1.09)\\
Fam. Income (75K to 100=1)   &0.06 (0.81)&0.01 (0.18)&-0.04 (-0.49)&-0.04 (-0.55)&0.13 (1.72)\\
Fam. Income ($>100K=1$)         &0.18 (2.21)&0.18 (2.27)&0.12 (1.53)&0.18 (2.20)&0.26 (3.16)\\
Gender (Men=1)                &-0.14 (-5.23)&-0.04 (-1.64)&-0.21 (-7.78)&-0.11 (-3.78)&-0.13 (-4.95)\\
Gender (Women=1)              &0.14 (5.23)&0.04 (1.64)&0.21 (7.78)&0.11 (3.78)&0.13 (4.95)\\
Race/Eth. (White=1)           &-0.02 (-0.34)&-0.11 (-1.53)&-0.00 (-0.03)&-0.18 (-2.45)&-0.16 (-2.18)\\
Race/Eth. (Black/Afr. Amer.=1)&-0.01 (-0.14)&0.02 (0.22)&0.16 (1.72)&-0.05 (-0.55)&0.14 (1.48)\\
Race/Eth. (Asian=1)           &0.07 (0.41)&-0.05 (-0.31)&-0.41 (-2.25)&-0.07 (-0.38)&0.17 (0.97)\\
Race/Eth. (Hispanic=1)        &-0.05 (-0.42)&-0.20 (-1.78)&-0.04 (-0.35)&-0.04 (-0.31)&-0.13 (-1.13)\\
Race/Eth. (Other=1)           &0.01 (0.07)&0.35 (2.05)&0.29 (1.67)&0.33 (1.99)&-0.02 (-0.13)\\ \\
Model $\chi^{2}$&308.17&   158.81&   179.81&   144.53&   176.80\\
N&     1702&     1706&     1704&     1702&     1704\\
\hline
\end{tabular*}
\label{tab2a}
\end{table}
%%\end{landscape}

%%\begin{landscape}
\begin{table}[t]
\caption{Coefficient estimates of the effects of socio-demographic covariates on rates of culture consumption obtained from a bivariate probit model with frequency of culture-talk and frequency of culture consumption within five domains as the joint as the joint outcome, 2002 and 2004 National Social Survey}
\begin{tabular*}{\hsize}{@{\extracolsep{\fill}} l*{5}{c}}
\hline
\hline
&\multicolumn{5}{c}{How often engage in:}\\
&Movies&Music&Books&Arts&Performances\\
\hline
Education (Less than H.S.=1)  &-0.36 (-4.16)&-0.07 (-0.75)&-0.41 (-4.89)&-0.44 (-4.58)&-0.41 (-4.68)\\
Education (High School=1)     &-0.05 (-1.07)&-0.06 (-1.10)&-0.22 (-4.44)&-0.38 (-6.71)&-0.22 (-4.36)\\
Education (Some Coll.=1)      &0.13 (2.59)&0.04 (0.64)&0.08 (1.62)&0.09 (1.71)&0.09 (1.82)\\
Education (College=1)         &0.09 (1.64)&-0.03 (-0.49)&0.18 (3.37)&0.20 (3.48)&0.16 (2.89)\\
Education (Postgraduate=1)    &0.20 (2.88)&0.12 (1.59)&0.36 (5.15)&0.53 (7.40)&0.39 (5.57)\\
Age (18-24=1)                 &0.73 (9.84)&0.59 (5.93)&0.06 (0.82)&0.15 (1.93)&0.35 (4.68)\\
Age (25-29=1)                 &0.38 (4.73)&0.14 (1.47)&-0.07 (-0.86)&0.09 (1.09)&0.17 (2.16)\\
Age (30-34=1)                 &-0.09 (-1.06)&0.07 (0.76)&-0.12 (-1.48)&0.02 (0.23)&-0.08 (-1.04)\\
Age (35-39=1)                 &0.10 (1.29)&0.09 (1.00)&-0.12 (-1.54)&-0.13 (-1.55)&-0.07 (-0.92)\\
Age (40-49=1)                 &-0.16 (-2.60)&-0.11 (-1.73)&-0.02 (-0.29)&-0.03 (-0.47)&-0.06 (-1.09)\\
Age (50-64=1)                 &-0.30 (-5.48)&-0.26 (-4.31)&0.07 (1.28)&-0.05 (-0.78)&-0.18 (-3.25)\\
Age (65+=1)                   &-0.67 (-10.24)&-0.51 (-7.83)&0.20 (3.17)&-0.06 (-0.92)&-0.12 (-1.91)\\
Fam. Income ($<20K=1$)          &-0.34 (-4.75)&0.00 (0.06)&0.01 (0.10)&-0.14 (-1.83)&-0.27 (-3.71)\\
Fam. Income (20K to 35K=1)   &-0.22 (-3.84)&-0.06 (-1.00)&-0.06 (-0.99)&-0.19 (-3.05)&-0.23 (-3.87)\\
Fam. Income (35K to 50K=1)   &-0.13 (-2.54)&-0.09 (-1.60)&-0.06 (-1.13)&-0.08 (-1.49)&-0.06 (-1.10)\\
Fam. Income (50K to 75K=1)   &0.13 (2.30)&0.08 (1.21)&0.02 (0.37)&0.06 (1.02)&0.17 (3.06)\\
Fam. Income (75K to 100=1)   &0.23 (2.98)&-0.01 (-0.15)&-0.02 (-0.22)&0.08 (1.07)&0.16 (2.12)\\
Fam. Income ($>100K=1$)         &0.34 (4.21)&0.09 (0.93)&0.10 (1.26)&0.27 (3.24)&0.22 (2.67)\\
Gender (Men=1)                &-0.07 (-2.43)&-0.07 (-2.49)&-0.16 (-6.12)&-0.05 (-1.60)&-0.01 (-0.50)\\
Gender (Women=1)              &0.07 (2.43)&0.07 (2.49)&0.16 (6.12)&0.05 (1.60)&0.01 (0.50)\\
Race/Eth. (White=1)           &-0.09 (-1.30)&-0.19 (-2.19)&-0.05 (-0.70)&-0.07 (-0.89)&0.06 (0.84)\\
Race/Eth. (Black/Afr. Amer.=1)&0.07 (0.70)&0.11 (0.96)&0.14 (1.50)&-0.04 (-0.36)&0.16 (1.68)\\
Race/Eth. (Asian=1)           &0.08 (0.47)&-0.02 (-0.11)&-0.45 (-2.51)&0.15 (0.85)&-0.28 (-1.54)\\
Race/Eth. (Hispanic=1)        &0.33 (2.82)&-0.25 (-1.86)&0.18 (1.56)&-0.02 (-0.19)&0.18 (1.59)\\
Race/Eth. (Other=1)           &-0.38 (-2.22)&0.36 (1.61)&0.18 (1.02)&-0.03 (-0.16)&-0.12 (-0.71)\\ \\
Model $\chi^{2}$                       &   308.17&   158.81&   179.81&   144.53&   176.80\\
N                             &     1702&     1706&     1704&     1702&     1704\\
\hline
\end{tabular*}
\label{tab2b}
\end{table}
%%\end{landscape}

%%\begin{landscape}
\begin{table}[t]
\caption{Coefficient estimates of the endogeneous effect ($\delta_{21}$) of frequency of consumption on culture-talk within five domains obtained from a recursive bivariate probit model, 2004 National Social Survey.}
\begin{tabular*}{\hsize}{@{\extracolsep{\fill}} l*{5}{c}}
\hline
\hline
&\multicolumn{5}{c}{Talk to someone you know about:}\\
&Movies&Music&Books&Arts&Drama\\
\hline
Frequency of:&&&&\\
Going to the Movies&0.54 (16.69)&&&\\
Music Listening&&0.34 (10.67)&&&\\
Reading&&&0.48 (18.63)&&\\
Arts Attendance&&&&0.76 (20.78)&\\
Performance Attendance&&&&&0.55 (17.77)\\
\hline
\end{tabular*}
\label{tab3}
\end{table}
%%\end{landscape}

%%\begin{landscape}
	\begin{table}[htb!]
	\subfloat [Coefficient estimates of the endogeneous effect ($\delta_{21}$) of different forms of musical taste obtained from a recursive bivariate semi-ordered probit model with frequency of culture-talk involving music as the joint outcome, 2004 National Social Survey. ] {
		\begin{tabular*}{\hsize}{@{\extracolsep{\fill}} l*{7}{c}} 
		\hline 
		\hline 
		&\multicolumn{6}{c}{Dependent Variable: Talk About Music with Friends}\\
		\hline
		Favorite Music?&&&&&&\\
		Opera, Classical=1&0.83 (13.77)&&&&&\\
		Jazz, Blues=1&&0.19 (3.78)&&&&\\
		Rock, Pop=1&&&0.11 (2.55)&&&\\
		Punk, Heavy Metal=1&&&&0.16 (1.70)&&\\
		Christian, Gospel=1&&&&&-0.03 (-0.63)&\\
		Folk, Country=1&&&&&&-0.67 (-8.64)\\ \\
		Less than H.S.&-0.05 (-0.48)&-0.28 (-2.41)&-0.29 (-2.57)&-0.38 (-3.47)&-0.32 (-2.80)&-0.10 (-0.84)\\
		High School&-0.11 (-1.50)&-0.23 (-2.99)&-0.26 (-3.33)&-0.26 (-3.46)&-0.24 (-3.14)&-0.04 (-0.55)\\
		Some Coll.&0.05 (0.63)&0.12 (1.58)&0.13 (1.74)&0.13 (1.74)&0.13 (1.74)&0.07 (1.02)\\
		Bachelor  &0.03 (0.40)&0.13 (1.60)&0.12 (1.54)&0.17 (2.22)&0.13 (1.69)&0.02 (0.21)\\
		Graduate  &0.09 (0.89)&0.26 (2.53)&0.30 (2.90)&0.35 (3.43)&0.30 (2.88)&0.05 (0.47)\\
		\hline
		
		\hline
	\end{tabular*}
	\label{tab4a}
	} \\
	\subfloat [Coefficient estimates of the effect of education on musical taste from a recursive bivariate semi-ordered probit model with frequency of culture-talk involving music as the joint outcome, 2004 National Social Survey.] {
		\begin{tabular*}{\hsize}{@{\extracolsep{\fill}} l*{7}{c}} 
		\hline 
		\hline 
		&\multicolumn{6}{c}{Dependent Variable: My favorite type music is:}\\
		&Opera, Classical&Jazz, Blues&Rock, Pop&Punk, Heavy Metal&Christian, Gospel&Folk, Country\\
		\hline
		Less than H.S.=1&-0.74 (-4.16)&-0.53 (-2.95)&-0.38 (-2.73)&0.21 (0.87)&0.17 (1.20)&0.36 (2.78)\\
		High School=1&-0.19 (-1.71)&-0.09 (-0.82)&0.23 (2.54)&0.16 (0.88)&-0.21 (-1.98)&0.30 (3.20)\\
		Some Coll.=1&0.23 (2.33)&0.14 (1.41)&-0.00 (-0.02)&0.13 (0.76)&0.04 (0.44)&-0.06 (-0.72)\\
		College=1&0.22 (2.02)&0.11 (1.02)&0.15 (1.60)&-0.05 (-0.27)&0.05 (0.47)&-0.19 (-1.94)\\
		Postgrad.=1 &0.48 (3.66)&0.37 (2.72)&-0.00 (-0.01)&-0.44 (-1.22)&-0.05 (-0.35)&-0.40 (-3.08)\\ \\
		N&\multicolumn{6}{c}{750}\\ 
		\hline 
	\end{tabular*}
	\label{tab4b}
	}	
	\caption{}
	\label{tab4}
\end{table}
%%\end{landscape}
\end{document}
